\vspace{-8pt}
\section{Language Generation through Multiagent Debate}
\vspace{-3pt}

We present an approach to generate language responses through multiagent debate. We provide an overview of our approach in \sect{sect:multiagent}. We further discuss convergence to consensus in the debate process in \sect{sect:consensus}. The overall overview of our approach is shown in \fig{fig:debate_overview}.

\begin{figure}[t]
    \centering
    \includegraphics[width=\linewidth]{fig/fig2-2.pdf}
    \caption{\textbf{Illustration of Debate.} Illustration of the debate procedure.}
    \label{fig:debate_overview}
    \vspace{-10pt}
\end{figure}



\subsection{Multiagent Language Generation}
\label{sect:multiagent}
\vspace{-3pt}

Consider your work process when solving the following math question on an exam: {\it ``What is the area of a triangle with side lengths of 3, 4, 5?"}. In one thread of work, you may recognize that the triangle side-lengths directly correspond to a right triangle, and thus directly compute the area as $0.5 \times 3 \times 4 = 64$.  To make sure that you have the right answer, you may then try to solve the problem differently by estimating an angle $\theta$ in the triangle using the Law of Cosines, and then obtain the area by using the formula $0.5 \times 3 \times 4  \times \sin(\theta)$, arriving at another answer to the given exam problem.  

When these lines of work give the same answer, your confidence about the answer increases. In contrast, when these answers are different, individual lines of work may engage in a mental ``debate" procedure, where you closely cross-examine the reasoning and assumptions of each line of work and refine solutions  until a consistent answer.

Similarly, consider writing a biography of a historical figure. To ensure the factuality of the biography, you may consult multiple different sources on each fact. Facts that are consistent in each source increase your confidence about the fact. In contrast, facts that are inconsistent require careful cross-examination between sources to determine the final consistent data.


To mimic the above multi-threaded 
 reasoning process and multi-source factuality checking processes, we propose to generate answers subject to a multi-agent debate procedure between multiple instances of large language models. Given a question, multiple agents represented as copies of a large language model, generate answers to the question. Each response serves as a possible thought process or source of information which agents may re-examine to find consistent final answers.

After initial responses are generated from different agents, we initiate a round of debate between agents. Individual responses from other agents are concatenated and given as context to each agent, with each agent instructed to construct a new response based on such responses. Each language agent is thus responsible for both verifying the collection of responses given by other agents, and refining its own response based on other agents' responses. We iteratively repeat this debate procedure over multiple rounds for improved performance. 

Concretely, we first prompt each agent to independently solve the given problem or task. After each agent generates a response, we feed each agent a consensus prompt, illustrated in \fig{tbl:prompt}, where each agent is instructed to update their responses based on the responses of other agents. This resultant consensus prompt may then be repeatedly given, using the updated responses of each agent.
We illustrate an overview of this multiagent debate procedure in \fig{fig:debate_overview}.


Note that our proposed approach operates in an orthogonal manner to existing approaches to prompt language models. Given a question, we may apply additional techniques for prompting language models to further improve our debate procedure by eliciting additional more detailed responses from language models. We illustrate the synergy of our approach with existing approaches to prompting language models in \fig{fig:prompt_ablation} and directly apply zero-shot chain-of-thought reasoning in our evaluations.


\begin{table*}[t]
\small\setlength{\tabcolsep}{5.5pt}
\centering
\scalebox{0.9}{

\begin{tabular}{c|c}
     {\bf Debate Length} & {\bf Prompt} \\
      \midrule
     \multirow{2}{*}{Short}  & \emph{" These are the solutions to the problem from other agents: [other answers]} \\
      & \emph{Based off the opinion of other agents, can you give an updated response $\ldots$"} \\
      \midrule
      \multirow{2}{*}{Long}  & \emph{" These are the solutions to the problem from other agents: [other answers]}  \\
     & \emph{Using the opinion of other agents as additional advice, can you give an updated response $\ldots$"} \\
    \bottomrule
\end{tabular}
}
\captionof{figure}{\small {\bf Prompts to induce long and short form debate.} Responses of other agents to questions are are inserted in the middle of the prompt (indicated with \emph{[other answers]})}
\label{tbl:prompt}
\vspace{-10pt}
\end{table*}

\vspace{-5pt}
\subsection{Consensus in Debates}
\label{sect:consensus}

Given multiple rounds of debate, how can we ensure that a set of  language model agents will converge to a final consensus answer? In general, debate can be seen as a multi-agent game, where convergence is not guaranteed. Empirically, however, we find that language models are able to converge on a single shared answer after multiple rounds of debate (\fig{fig:math_overview}).

We found that we could control the duration of debates by how changing how much a language model trusts its own outputs over those generated by other models through different prompts. We illustrate two prompts below in \fig{tbl:prompt}, which we use to induce different debate durations between language models, and illustrate the effect of such prompts in \fig{fig:convergence}. In general, we found that prompts that encouraged models to be more ``stubborn’ based on their own solutions led to longer debates and better final solutions. Overall, we observed that language model agents were relatively "agreeable", perhaps as a result of instruction tuning or reinforcement learning based on human feedback ~\cite{ouyang2022training}.


